\documentclass[11pt,a4paper]{article}
\usepackage[margin=1in]{geometry}
\usepackage{amsmath,amssymb,amsthm}
\usepackage{booktabs}
\usepackage{xcolor}
\definecolor{linkcol}{rgb}{0.05,0.3,0.45}
\usepackage[colorlinks=true,linkcolor=linkcol,citecolor=linkcol,urlcolor=linkcol]{hyperref}

\newtheorem{theorem}{Theorem}
\newtheorem{proposition}{Proposition}
\newtheorem{lemma}{Lemma}
\newtheorem{corollary}{Corollary}
\newtheorem{observation}{Observation}
\theoremstyle{definition}
\newtheorem{example}{Example}
\newtheorem{remark}{Remark}

\newcommand{\ket}[1]{\lvert #1 \rangle}
\newcommand{\bra}[1]{\langle #1 \rvert}
\newcommand{\proj}[1]{\lvert #1 \rangle\!\langle #1 \rvert}
\newcommand{\tr}{\operatorname{Tr}}
\newcommand{\id}{\mathbb{1}}
\newcommand{\F}{\mathcal{F}}
\newcommand{\Ops}{\mathcal{O}}

\title{The Quantum Resource Network: a role-resolved theorem graph of
thirteen resource theories with via-annotated conversions and
machine-visible open problems}

% TODO (author): fill in name(s), affiliation and email before submission.
% Consider disclosing AI-assisted drafting/curation if required by the venue's policy;
% the curation and verification protocol is described in Sec. 6.
\author{[Author name]\\ \small [Affiliation] \quad [email]}
\date{October 2026 \quad (v1.1 of the network)}

\begin{document}
\maketitle

\begin{abstract}
We present a curated, typed, signed graph of thirteen quantum resource
theories --- purity, athermality, asymmetry, coherence, imaginarity, magic,
continuous-variable Wigner negativity, entanglement, steering, Bell
nonlocality, contextuality, measurement incompatibility, and quantum discord
as the documented counterexample to the framework axioms --- released as an
open repository with an interactive page and machine-readable data. The design
resolves each theory into \emph{role-typed} nodes: free states versus resource
states; free operations versus \emph{resource-utilizing} operations versus
\emph{impossible-even-with-resources} processes (no-go theorems as first-class
nodes); and detector versus quantifier witnesses, a split made load-bearing by
the fact that facet Bell-inequality violations detect nonclassicality but do
not determine the convertibility pre-order. Cross-theory edges therefore land
on the exact object they concern, and the compilation enforces one convention
throughout: \emph{a conversion between resource theories is well-defined only
relative to an operation class}, so every conversion edge carries a
\texttt{via} field naming that class. Absence of an edge always means
``unsettled'', never ``false'', which turns the $13\times13$ directed
conversion matrix (13 of 156 cells charted at this release) into a
machine-generated research agenda. Alongside the compiled backbone --- the
free-set geography, the operations lattice, the 2023--2024 reversibility
frontier (irreversibility under exactly resource-non-generating operations
coexisting with the generalized quantum Stein's lemma), and the
necessity-pair structure of steering --- we record five derived propositions
obtained by closing gaps the compilation exposed, including a typology of
free-operation justifications whose unique two-type intersection
(athermality) reappears at the state and witness levels, the
operation-class-relativity of bound resources, and the observation that the
Gibbs state of a generic interacting Hamiltonian is a magic resource while
free in four other theories. Ten named open problems are registered. All
sources, data and pages are provided for community extension.
\end{abstract}

\section{Introduction}

Quantum resource theories share one skeleton --- a free set $\F$ and a class
of free operations $\Ops$ with the golden rule $\Ops(\F)\subseteq\F$
\cite{chitambargour2019} --- and the skeleton has produced a genuinely
transversal body of results: faithful conversion of coherence into
entanglement \cite{streltsov2015}, activation of discord into distillable
entanglement \cite{piani2011}, the one-to-one correspondence between steering
and measurement incompatibility \cite{uola2015}, the identification of purity
as the maximal unitarily reachable coherence \cite{streltsov2018njp}, and,
recently, a second law of resource manipulation restored at the level of
asymptotically resource-non-generating operations by the proof of the
generalized quantum Stein's lemma \cite{hayashiyamasaki2024,lami2024}
coexisting with its impossibility at the exact level
\cite{lamiregula2023}. Yet these results are stated across subcommunities in
different operation classes, and the most common source of confusion in the
cross-theory literature is precisely that ``a resource of $A$ yields a
resource of $B$'' is not a well-formed claim until the operation class is
fixed. A second difficulty is structural: survey prose does not distinguish
``no theorem connects $X$ to $Y$'' from ``a theorem forbids it''.

This paper introduces a working tool addressed at both difficulties: the
\emph{Quantum Resource Network}, a signed theorem graph over thirteen
resource theories, companion to the Nonclassicality Network of quantum
foundations released earlier in the same series. Four design decisions do
the work.

\begin{enumerate}
  \item \textbf{Role-resolved bands.} Each theory is a band of typed nodes:
  its model axioms; its \emph{free states} and its \emph{resource states} as
  separate nodes; its \emph{free operations}, its \emph{resource-utilizing
  operations} (what spending the resource enables: injection, activation,
  work extraction, discrimination advantages), and --- where a genuine
  theorem exists --- its \emph{impossible processes}, no-gos that hold even
  with unlimited resource (four at this release: thermal coherence
  distillation and work locking \cite{marvian2020,lostaglio2015natcomm};
  mixed-state irreversibility under incoherent operations
  \cite{winteryang2016}; the absence of a second law of entanglement under
  exactly non-entangling operations together with bound entanglement
  \cite{lamiregula2023,horodecki2009}; bound magic and the failure of
  unit-rate interconversion \cite{veitch2012}). Witnesses are split into
  \emph{detectors} and \emph{quantifiers}; the split is not pedantry, since
  for Bell nonlocality the degrees of violation of all facet inequalities
  detect but provably do not quantify \cite{wolfe2020}.

  \item \textbf{The via convention.} Every cross-theory conversion edge
  carries a \texttt{via} field naming the operation class under which the
  conversion holds (bipartite incoherent operations for
  coherence$\to$entanglement \cite{streltsov2015}; local
  quantum-incoherent operations and classical communication for the reverse
  direction \cite{chitambarhsieh2016}; local ancilla interactions plus
  distillation for discord activation \cite{piani2011,streltsov2011}; and so
  on). Conversions run resource node to resource node; free-set relations run
  free node to free node.

  \item \textbf{Three-valued edges.} Positive edges record possibility and
  structure theorems; negative edges record no-gos, strict separations and
  counterexamples, drawn with inhibition-style flat-bar terminators
  ($\dashv$); no edge means \emph{unsettled}. Without first-class negative
  edges, impossibility theorems would be encoded as absences and a search for
  open problems would re-propose closed questions.

  \item \textbf{Per-reference provenance.} Every reference string in the data
  carries one of three tags: \texttt{verified} (bibliographic data confirmed
  against sources on 2026-10-23 --- the load-bearing cross-theory anchors),
  \texttt{knowledge} (standard results cited from curation knowledge, page
  data to be rechecked before quotation), or \texttt{verify} (flagged
  uncertain details that must be checked before any release).
\end{enumerate}

\paragraph{Motivation and intended use.}
The network was built for three purposes. \emph{(i)~Automating research-idea
generation.} Because no edge means ``unsettled'', the open questions form a
queryable object. The canonical instance here is the $13\times13$ directed
conversion matrix, each cell the question ``can a resource of theory $A$
produce a resource of theory $B$, and under which operation class?'': this
release charts 13 of 156 directed cells, and the remaining 143, together with
the ten named problems of Sec.~\ref{sec:open}, constitute the generated
agenda. \emph{(ii)~Reviewing a field as curation.} Each claim of the review
is one edge with a verbatim statement, references and per-reference
provenance; checking the review means auditing edges, and disputing it means
objecting to an edge identifier. \emph{(iii)~Visualization.} The
role-resolved bands make visible at a glance where free sets overlap, where
resources convert and through which operations, where spending is powerful
and where it is provably futile, and which witnesses merely detect. The
review process is itself the generator: surveying the field in this format
yields its own research agenda as a by-product.

Section~\ref{sec:theories} fixes the thirteen theories.
Section~\ref{sec:backbone} summarizes the curated backbone.
Section~\ref{sec:derived} contains the derived propositions with proofs or
proof sketches. Section~\ref{sec:open} lists the registered open problems,
Section~\ref{sec:protocol} the curation and verification protocol, and
Section~\ref{sec:repo} the repository.

\section{The thirteen theories}\label{sec:theories}

Table~\ref{tab:theories} lists the compiled theories with their free sets,
physically motivated free-operation classes and principal monotones; the data
file carries per-row references. Twelve are proper resource theories; quantum
discord is carried as the documented counterexample: its classical--classical
free set is non-convex, and local channels can create discord
\cite{modi2012,ciccarello2012}, so no golden-rule operation class containing
local processing exists. Measurement incompatibility is carried as the
demonstration that the schema extends from states to devices: free objects
are classically programmable measurement devices, and convertibility is
completely characterized by guessing games with post-measurement information
\cite{bcz2020}.

\begin{table}[t]
\centering\small
\begin{tabular}{@{}lllll@{}}
\toprule
theory & type & free states $\F$ & free ops $\Ops$ & refs \\
\midrule
purity & constraint & $\{\id/d\}$ & noisy/unital & \cite{gour2015} \\
athermality & constraint$\times$symmetry & Gibbs $\gamma$ & thermal ops (TO) & \cite{brandao2013,horodeckioppenheim2013} \\
asymmetry & symmetry & $G$-invariant & $G$-covariant & \cite{gourspekkens2008,marvianspekkens2016} \\
coherence & dephasing-defined & diagonal & SIO/IO/DIO/MIO & \cite{baumgratz2014,streltsov2017rmp} \\
imaginarity & antiunitary $\mathbb{Z}_2$ & real states & real-Kraus ops & \cite{hickeygour2018,wu2021prl} \\
magic & subtheory & STAB & stabilizer/CSPO & \cite{veitch2014,seddoncampbell2019} \\
Wigner neg.\ (CV) & subtheory & Wigner-positive & Gaussian protocols & \cite{takagizhuang2018,marieisert2012} \\
entanglement & partition & SEP & LOCC & \cite{horodecki2009} \\
steering & partition (1-sided) & LHS assemblages & 1W-LOCC & \cite{wjd2007,gallegoaolita2015} \\
nonlocality & partition (no comm.) & local boxes & LOSR wirings & \cite{wolfe2020} \\
contextuality & partition (compat.) & noncontextual & NC wirings & \cite{amaral2018} \\
incompatibility & device & jointly measurable & pre/post-processing & \cite{bcz2020,heinosaari2016} \\
discord & (fails axioms) & CC (non-convex) & none & \cite{modi2012,ollivierzurek2001} \\
\bottomrule
\end{tabular}
\caption{The thirteen compiled theories. ``Type'' refers to the typology of
Proposition~\ref{prop:qp1}.}
\label{tab:theories}
\end{table}

\section{The curated backbone}\label{sec:backbone}

\paragraph{Free-set geography.} The free sets are pairwise transversal rather
than nested, and the network records the exceptions as explicit subset edges:
$\{\gamma\}\subset$ incoherent(energy basis) $\subset$ real(energy basis);
incoherent(product basis) $\subset$ SEP; the strict chain SEP $\subsetneq$
unsteerable $\subsetneq$ Bell-local
\cite{wjd2007,werner1989,barrett2002,quintino2015}, with Bell-local behaviors
embedding scenario-by-scenario into the noncontextual behaviors
\cite{amaral2018}; and, as a two-sided negative edge, the transversality of
STAB and SEP --- the Bell state is a stabilizer state, and
$\ket{T}\otimes\ket{0}$ is separable and magic. Distinguished shared states
carry the membership structure: $\id/d$ (free everywhere except purity),
the Gibbs state (free in four theories, generically a magic resource,
Proposition~\ref{prop:qp5}), $\ket{+}$ (maximal coherence, magic- and
imaginarity-free), the $T$ state (a four-theory resource
\cite{bravyikitaev2005}), Werner states (the strictness witness of the
hierarchy \cite{werner1989,barrett2002,wjd2007}), bound entangled states
(nonlocal despite $E_D=0$ \cite{vertesi2014}), and Wigner-positive magic
states (bound magic \cite{veitch2012}).

\paragraph{Operations lattice.} The partition ladder LOSR $\subset$ 1W-LOCC
$\subset$ LOCC $\subset$ SEP-ops $\subset$ PPT-ops $\subset$ NE generates
nonlocality, steering and entanglement from one bipartite scenario
\cite{gallegoaolita2015,wolfe2020}; LOCC is the wrong class for Bell
scenarios, where the relevant resource is LOSR-entanglement \cite{wolfe2020}.
On the symmetry side, TIO $\subset$ DIO and the coherence
lattice SIO $\subset$ IO, DIO $\subset$ MIO \cite{marvianspekkens2016};
thermal operations are simultaneously Gibbs-preserving and covariant, and
strictly inside the Gibbs-preserving maps because the latter can create
energy coherence \cite{faist2015}. Noncontextual wirings restricted to Bell
scenarios differ from LOSR \cite{santos2024}, so the two candidate
unifications of ``classical operations'' disagree (QG3).

\paragraph{Conversions, with their via classes.} Coherence converts to
entanglement faithfully under bipartite incoherent operations, with
$E_r \le C_r$ attainable \cite{streltsov2015}; entanglement converts back to
coherence in the intersection class LQICC \cite{chitambarhsieh2016}; purity
is the maximal unitarily reachable coherence and bounds unitary entanglement
generation \cite{streltsov2018njp}; every discordant state --- separable ones
included --- activates into distillable entanglement, with the relative
entropy of quantumness equal to the minimal entanglement generated
\cite{piani2011,streltsov2011}; entanglement yields steering through
incompatible measurements on the untrusted wing
\cite{wjd2007,quintino2014,uola2015} and nonlocality through local
measurements, insufficiently per copy \cite{werner1989,barrett2002} but
superactivatably \cite{palazuelos2012}, with every entangled state winning
some semiquantum game \cite{buscemi2012}; magic yields contextuality under
stabilizer measurements, with equivalence to distillability in odd prime
dimension \cite{howard2014} and a continuous-variable twin equivalence
\cite{booth2022,marieisert2012}; imaginarity buys a local-discrimination gap
that LOCC with real operations provably cannot close \cite{wu2021prl,wu2021pra}.
Catalysis is carried as a shared operations-layer phenomenon
\cite{jonathanplenio1999,aberg2014,datta2023}.

\paragraph{The reversibility frontier.} Under \emph{exactly} non-entangling
operations, entanglement manipulation is asymptotically irreversible and no
single measure governs all transformations \cite{lamiregula2023}; under
\emph{asymptotically} resource-non-generating operations, every theory
satisfying the Brand\~ao--Plenio axioms is reversible with the regularized
relative entropy of resource as the unique rate function --- the generalized
quantum Stein's lemma, claimed in 2010, whose proof gap was identified in
2023 \cite{berta2023} and closed in 2024
\cite{hayashiyamasaki2024,lami2024,brandaogour2015}. The exact/asymptotic
distinction is therefore the load-bearing hinge of second-law statements
(QG4). The universal witness layer is organized by the robustness theorems:
in any convex theory every resource state strictly outperforms all free
states in some subchannel discrimination task, with the generalized
robustness as the exact advantage
\cite{takagi2019,takagiregula2019,vidaltarrach1999}, extended beyond
convexity in multi-copy form \cite{kuroiwa2024}; the steering instance is
discrimination under one-way LOCC measurements \cite{pianiwatrous2015}.

\section{Derived propositions}\label{sec:derived}

The following were assembled by closing gaps the compilation exposed. They
are proved or sketched here and in the repository catalog; \emph{no novelty
pass has been run} (unlike the companion repository's generated
propositions), QP2--QP4 are plausibly folklore, and none is claimed as a
first. Statements are numbered as in the data files.

\begin{proposition}[QP1: typology of free-operation justifications]
\label{prop:qp1}
Within the compiled set, every proper theory's free operations are justified
in exactly one of four ways --- symmetry (covariance under a group action:
asymmetry; unspeakable coherence; imaginarity as the antiunitary
$\mathbb{Z}_2$ analogue), subtheory (operations of a classically simulable
subtheory: magic, CV non-Gaussianity), partition (compatibility with a causal
or communication partition: entanglement, steering, nonlocality,
contextuality, incompatibility), or constraint (preservation of a
distinguished state: purity) --- with a single two-type intersection:
athermality, whose thermal operations are simultaneously Gibbs-preserving
(constraint) and time-translation covariant (symmetry). The intersection
reappears at the state level, where the free set is the unique state fixed by
both axioms (the Gibbs state), and at the witness level, where the complete
second-law constraints factor into the R\'enyi free-energy family on
block-diagonal data \cite{brandao2015pnas} and independent asymmetry-mode
constraints on coherences \cite{lostaglio2015prx}, neither family alone
complete \cite{lostaglio2015natcomm}; the two parent operation classes
separate strictly in the quantum regime \cite{faist2015}.
\end{proposition}

\begin{proof}[Proof sketch]
The classification is verified theory by theory against the defining
free-operation classes of Table~\ref{tab:theories}. That athermality
instantiates both axioms is the definition of thermal operations
(energy-conserving unitaries with Gibbs ancillas); uniqueness of the
two-axiom fixed point for nondegenerate $H$ follows since a U(1)-invariant
state is energy-diagonal and a Gibbs-preserving fixed point with full rank is
$\gamma$. The witness-level factorization is the content of
\cite{brandao2015pnas,lostaglio2015prx,lostaglio2015natcomm}.
\end{proof}

\begin{observation}[QP2: tensor-stability failure of state-level Bell
locality]\label{obs:qp2}
The set of Bell-local quantum states is not closed under tensor products:
there exist states admitting local-hidden-variable models for all single-copy
measurements whose tensor powers violate a Bell inequality
\cite{palazuelos2012}. Hence ``Bell-local states'' fail the monoidal
free-set axiom of the resource-theory framework \cite{cfs2016}, while
common-cause \emph{boxes} under LOSR wirings satisfy it \cite{wolfe2020}.
The diagnosis: when a framework axiom fails, the objects were wrongly chosen,
not the axiom.
\end{observation}

\begin{observation}[QP3: bound resources are operation-class-relative]
\label{obs:qp3}
Whether a theory exhibits bound (undistillable) resources is a joint property
of the pair $(\F,\Ops)$, not of the free set alone. Documented cases:
entanglement under LOCC has bound states (PPT entangled, $E_D=0<E_C$)
\cite{horodecki2009}; magic under stabilizer operations has bound states
(Wigner-positive non-stabilizer states) \cite{veitch2012}; coherence under
thermal/covariant processing is \emph{entirely} bound at full rank ---
distillable coherence vanishes while the cost equals the quantum Fisher
information \cite{marvian2020}; coherence under incoherent operations has
\emph{no} bound states --- every coherent state distills at rate $C_r$
\cite{winteryang2016}; purity under noisy operations has none
\cite{gour2015}. The coherence pair is the sharp instance: one free set,
two operation classes, opposite phenomenology.
\end{observation}

\begin{proposition}[QP4: the purity--coherence--entanglement cascade]
\label{prop:qp4}
For the relative-entropy quantifiers, the purity of a state equals its
maximal coherence reachable by free unitaries,
$P_r(\rho)=\max_U C_r(U\rho U^\dagger)$, and this maximal coherence bounds
the entanglement generable from the state by bipartite incoherent operations
acting on the state and an incoherent ancilla,
$E_r(\Lambda_{\mathrm{IO}}[\rho\otimes\sigma_{\mathrm{inc}}])\le C_r(\rho)$,
the bound attainable by a generalized CNOT. Thus
$P_r \ge C_r^{\max} \ge E_r^{\max}$, a quantitative generation chain across
three theories with each link's operation class stated.
\end{proposition}
\begin{proof}
The first equality is Theorem content of \cite{streltsov2018njp} (which also
shows purity bounds unitary entanglement and discord generation); the second
inequality and its attainability are the main results of
\cite{streltsov2015}.
\end{proof}

\begin{proposition}[QP5: thermal magic]\label{prop:qp5}
Let $H$ be a Hamiltonian on a finite-dimensional multi-qudit system whose
ground state $\ket{g}$ (assumed unique) is not a stabilizer state. Then there
exists $\beta_0<\infty$ such that for all $\beta>\beta_0$ the Gibbs state
$\gamma(\beta)=e^{-\beta H}/Z$ lies outside the stabilizer polytope: the
unique free state of athermality --- simultaneously free for coherence,
asymmetry and imaginarity in the energy eigenbasis --- is a magic resource.
\end{proposition}
\begin{proof}
The stabilizer polytope STAB is closed (a convex hull of finitely many
states). As $\beta\to\infty$, $\gamma(\beta)\to\proj{g}$ in trace norm, and
$\proj{g}\notin$ STAB by hypothesis; since the complement of STAB is open,
$\gamma(\beta)\notin$ STAB for all $\beta$ beyond some $\beta_0$. The curve
starts inside ($\gamma(0)=\id/d\in$ STAB), so a boundary crossing
$\beta^*\le\beta_0$ exists. Membership of $\gamma$ in the other three free
sets is elementary for nondegenerate $H$: $[\gamma,H]=0$ gives energy-basis
diagonality (incoherence and U(1)-invariance), and the diagonal entries are
real.
\end{proof}

\begin{remark}
The hypothesis holds for generic interacting Hamiltonians (non-stabilizer
ground states are the rule; stabilizer ground states are the exceptional,
e.g.\ Clifford-diagonalizable, case), but ``generic'' is used here
informally; the proposition as stated is the rigorous conditional. The
quantitative program --- mana or robustness of magic of $\gamma(\beta)$
across model families, and its relation to the classical simulability of
thermal states --- is registered as gap QG5. We did not locate a dedicated
treatment in the curation pass; a literature search before any claim of
novelty is mandatory (Sec.~\ref{sec:protocol}).
\end{remark}

\section{Registered open problems}\label{sec:open}

QG1: the exact scope of the multiqubit (even-$d$) magic--contextuality
correspondence \cite{howard2014}. QG2: is TO $=$ GPO $\cap$ TIO exactly? The
strict inclusions are settled \cite{faist2015}; the characterization is not.
QG3: the correct single free-operation class unifying contextuality and
nonlocality, given NCW$|_{\mathrm{Bell}}\neq$ LOSR \cite{santos2024,amaral2018}.
QG4: the reversibility-frontier table --- for each theory, reversibility
under exactly-RNG, asymptotically-RNG \cite{hayashiyamasaki2024,lami2024},
and the physically motivated class; known points: entanglement (no/yes/no)
\cite{lamiregula2023}, IO-coherence (no for cost/yes/partial)
\cite{winteryang2016,brandaogour2015}, thermal asymmetry (no, maximally)
\cite{marvian2020}, purity (yes/yes/yes) \cite{gour2015}. QG5: the uncharted
cells of the $13\times13$ via-class conversion matrix (thermal magic; real
thermal operations; the imaginarity cost of $T$ states), the automation hook
of the network. QG6: a steering analogue of the entangleability of cones.
QG7: an erasure-cost calculus beyond entanglement. QG8: a maximal physically
motivated operation class for non-convex free sets (the discord repair),
beyond the advantage result of \cite{kuroiwa2024}. QG9: catalytic and
battery routes around exact-RNG irreversibility \cite{datta2023}. QG10: the
thermodynamics of bound entanglement and of bound magic.

\section{Curation and verification protocol}\label{sec:protocol}

Reliability is handled by explicit provenance rather than by implication.
(i)~Every edge carries a verbatim statement and references; the per-edge
reference list, not this paper, is the citation of record for compiled
results. (ii)~References are individually tagged: \texttt{verified} entries
had their bibliographic data confirmed against sources on 2026-10-23
(approximately twenty load-bearing anchors, including
\cite{streltsov2015,marvianspekkens2016,msz2016,faist2015,lostaglio2015prx,
marvian2020,lamiregula2023,hayashiyamasaki2024,lami2024,wolfe2020,uola2015,
amaral2018,santos2024,wu2021prl,wu2021pra,bcz2020,takagi2019,takagiregula2019,
piani2011,streltsov2018njp,datta2023}); \texttt{knowledge} entries are
standard results cited from curation knowledge whose page data should be
rechecked before quotation; \texttt{verify} entries flag uncertain author
lists, venues or page numbers that must be checked before any release.
(iii)~Unlike the companion repository, the bibliography of this paper has
\emph{not yet} been verified programmatically against Crossref, and no
novelty searches have been run for QP1--QP5; both passes are advised before
submission. QP2--QP4 are assemblies of published results whose framing may
well exist in the literature; QP5's conditional statement is elementary.
(iv)~Derived items state exactly which recorded results they assemble, so a
reader can audit each in one step. (v)~The authors take responsibility for
remaining errors and welcome corrections as repository issues; the graph is
versioned, and every release records its census (v1.1: 116 nodes; 85 edges
--- 38 possibility, 13 no-go, 34 structural; 13/156 conversion cells
charted).

\section{Repository}\label{sec:repo}

The repository (\url{https://github.com/tahmeedhossain/GEN_QRT_Level},
served as live pages at
\url{https://tahmeedhossain.github.io/GEN_QRT_Level/}) contains:
\texttt{index.html} (overview, derived propositions, gap census);
\texttt{qrt.html} (the interactive network: 13 role-resolved theory bands,
edge catalog, theory table); \texttt{data/qrt\_network.json}
(machine-readable: theories, role-typed nodes, via-annotated signed edges
with statements and tagged references, derived propositions, gaps);
\texttt{catalog/qrt\_catalog.md} (prose catalog with the axiom-overlap
analysis, the free-set membership table and proof sketches); and this
paper's source under \texttt{paper/}. The data schema is documented in the
README; contributed edges are expected to carry a statement, a reference or
proof, a kind, a polarity, a status and --- for conversions --- the
\texttt{via} operation class. Closing a cell of the conversion matrix is the
canonical contribution.

\paragraph{Acknowledgements.} The curation of the network, the derivations
and drafting of this paper, the interactive pages and the data files were
produced with substantial assistance from an AI system (Claude, Anthropic),
working under the direction of the author, who reviewed the material and
takes sole responsibility for all claims and any remaining errors.
% TODO (author): add human acknowledgements/funding if appropriate.

\bibliographystyle{unsrt}
\bibliography{refs}

\end{document}
